% !TEX root = ../../tree_notes.tex
% =====================================================================
%  3.1  The Three Depth-First Orders                        sec:c03-dfs
%  Exports: def:c03-order  def:c03-three-orders  thm:c03-one-walk
%           thm:c03-root-position  cor:c03-rev-post  lst:c03-dfs
%           ex:c03-runs  prob:c03-preorder  prob:c03-inorder
%           prob:c03-verify-pre
% =====================================================================
\section{The Three Depth-First Orders}\label{sec:c03-dfs}

\begin{topicintro}
Chapter~2 recursed over a tree but never \emph{named} the moment a node's value was
used. This section names it. A \term{depth-first order} is one of three fixed
sequences in which a single recursive walk can \emph{emit} its nodes --
\term{preorder} (node, left, right), \term{inorder} (left, node, right) and
\term{postorder} (left, right, node). They are the same walk with the output
statement moved to one of three positions, and that one-line difference is the whole
reason three orders exist: each hands you a \emph{different set of already-known
facts} at the instant you act.
\textbf{Prerequisites.} \secref{c02-model} (\code{TreeNode}, structural recursion, the
base-case discipline); the running example $R$ (\figref{c03-running}).

\textbf{After this section you can}
\begin{itemize}[nosep]
  \item produce all three DFS orders recursively from one nine-line template;   \hfill\lv{1}
  \item read a problem's phrasing and know \emph{which} emit-point it is asking for; \hfill\lv{2}
  \item state exactly what is already computed when a node is visited in each order; \hfill\lv{2}
  \item decide, without building anything, whether a token string is a legal preorder
        of \emph{some} binary tree;                                             \hfill\lv{4}
\end{itemize}
\end{topicintro}

\begin{recognitionbox}[Pattern card -- a fixed depth-first visit order]
\reckey{``return the preorder / inorder / postorder traversal'', ``visit the root
before / between / after its subtrees'' -- the problem is defined by \emph{the order}
it wants nodes emitted.}
\recuse{
  \item the required \emph{output sequence} is the answer itself (print/store the
        traversal);
  \item you must act when a node is \emph{first seen} (preorder: copy a tree, encode a
        root) or only after \emph{both subtrees are done} (postorder: delete, aggregate);
  \item you need left-before-right structure explicitly (inorder: the sorted read of a
        BST, the left-edge of a subtree).
}
\recnotuse{
  \item order genuinely does not matter (bare subtree sums) -- the plain Chapter~2
        recursion (\secref{c02-bottomup}) already suffices;
  \item the answer is per \emph{row} / shortest-distance -- that is level order,
        \secref{c03-bfs};
  \item recursion is forbidden or extra space must be $\bigO{1}$ -- the stack/Morris
        forms live in \secref{c03-iterative};
  \item you are \emph{given} orders and must recover the tree -- reconstruction is
        \secref{c03-rebuild}.
}
One-line test: \emph{does the problem change if I permute the visit sequence? If yes,
you are choosing an emit-point -- this section.}
\end{recognitionbox}

\subsection{Intuition}

Run a pointer around the \emph{outside} of the tree, starting under the root and
keeping the left edge touching, exactly as if you were tracing the silhouette of the
drawing with your finger. That single contour pass visits every node \textbf{three
times} -- at its left shoulder, directly beneath it, and at its right shoulder. The
three depth-first orders are nothing but the same walk read at three different
shoulders:

\begin{figure}[H]
\centering
\begin{tikzpicture}[diagram, scale=0.85, transform shape]
  \node[troot] (r) at (0,0) {$1$};
  \node[tnode] (a) at (-1.7,-1.4) {$2$};
  \node[tnode] (b) at (1.7,-1.4) {$3$};
  \node[tleaf] (c) at (-2.6,-2.8) {$4$};
  \node[tleaf] (d) at (-0.8,-2.8) {$5$};
  \node[tnode] (e) at (1.7,-2.8) {$6$};
  \draw[tedge] (r) -- (a); \draw[tedge] (r) -- (b);
  \draw[tedge] (a) -- (c); \draw[tedge] (a) -- (d);
  \draw[tedge] (b) -- (e);
  % three markers under each node: pre (blue) / in (orange) / post (green)
  \foreach \x/\y in {0/0,-1.7/-1.4,1.7/-1.4,-2.6/-2.8,-0.8/-2.8,1.7/-2.8} {
    \node[font=\scriptsize,text=primaryblue]  at (\x-0.30,\y-0.48) {$\scriptstyle 1$};
    \node[font=\scriptsize,text=accentorange] at (\x      ,\y-0.55) {$\scriptstyle 2$};
    \node[font=\scriptsize,text=successgreen] at (\x+0.30,\y-0.48) {$\scriptstyle 3$};
  }
  \node[annot, anchor=west, align=left, font=\scriptsize] at (3.1,-0.4)
    {\textcolor{primaryblue}{$1$}=left shoulder\\(preorder)\\
     \textcolor{accentorange}{$2$}=underneath\\(inorder)\\
     \textcolor{successgreen}{$3$}=right shoulder\\(postorder)};
\end{tikzpicture}
\caption{One contour walk, three read-points. Passing the \emph{left} shoulder of a
node emits it in preorder; passing \emph{underneath} emits it in inorder; passing the
\emph{right} shoulder emits it in postorder. Read the numbers left-to-right per colour
to recover each traversal of $R$.}\label{fig:c03-contour}
\end{figure}

The picture explains why the three are siblings: they share the walk, so they share the
$\bigO{n}$ cost, the $\bigO{h}$ stack and the base-case discipline of
\lemref{c02-visit-count}. They differ only in \emph{what you know about the node at the
moment you emit it}. At its left shoulder you know nothing below it yet (preorder:
perfect for handing a value \emph{down}); underneath you know the entire left subtree
(inorder: the sorted read); at its right shoulder you know both subtrees completely
(postorder: the aggregate-and-delete read). That sentence is the decision table for
every problem in this chapter.

\subsection{Definitions and Notation}

\begin{definition}{Traversal (order)}{c03-order}
A \term{traversal} of a rooted tree is a rule that lists every node exactly once, in a
sequence determined by the tree's structure. It is \term{depth-first} if, once it
descends into a child subtree, it finishes that entire subtree before touching the next
sibling; it is \term{breadth-first} (\secref{c03-bfs}) if it exhausts one depth before
any node deeper.
\end{definition}

\noindent\textbf{In plain words.} A traversal is a \emph{route}, not new machinery: the
nodes are fixed, only the order you write them down changes. ``Depth-first'' means
``fully commit to one branch before backtracking'' -- which is exactly what a recursive
call does by nature.

\begin{definition}{The three depth-first orders}{c03-three-orders}
For a node $v$ with left subtree $L$ and right subtree $R$, the order of $\subtree(v)$ is
defined recursively by \emph{where the visit of $v$ sits} relative to the two sub-walks:
\[
\mathit{pre}(v) = v,\ \mathit{pre}(L),\ \mathit{pre}(R)\qquad
\mathit{in}(v)  = \mathit{in}(L),\ v,\ \mathit{in}(R)\qquad
\mathit{post}(v)= \mathit{post}(L),\ \mathit{post}(R),\ v.
\]
A missing subtree contributes nothing (the $\nil$ base case emits an empty list).
\end{definition}

\noindent\textbf{In plain words.} ``Pre'' visits the node \emph{before} its children,
``post'' \emph{after}, ``in'' \emph{between} (left first, then the node, then right).
The names describe the node's position in its own subtree's list -- nothing more.

\begin{vsbox}{Where $v$ sits in $\mathit{pre}/\mathit{in}/\mathit{post}$}{preorder}{in/post}
{ $v$ is \textbf{first}; the whole subtree follows it. Use when acting on a node must
happen \emph{before} its children exist/change. }
{ in inorder $v$ is \textbf{in the middle} (splits $L\mid R$); in postorder $v$ is
\textbf{last} (both subtrees precede it). Use inorder to separate the two children,
postorder to finish children first. }
\end{vsbox}

\subsection{Properties}

\begin{theorem}{One Walk, Three Outputs}{c03-one-walk}
A single depth-first recursion emits any one of the three orders depending only on the
position of the output statement relative to the two recursive calls. All three visit
each node exactly once and cost $\bigO{n}$ time and $\bigO{h}$ stack.
\end{theorem}
\begin{proofbox}
Fix the procedure \code{walk}(v): (i) if $v=\nil$ return; (ii) \code{emit\_pre};
\code{walk}(left); (iii) \code{emit\_in}; \code{walk}(right); (iv) \code{emit\_post}.
This is structural recursion on $\subtree(v)$ (rule (ii) of \thmref{c02-recursive-shape}),
so by \lemref{c02-visit-count} it makes $2n+1$ calls and visits every real node exactly
once, contributing it to exactly one of the three streams at each of the three marked
points of \figref{c03-contour}. Recording \code{emit\_pre} (resp.\ \code{emit\_in},
\code{emit\_post}) as the output reproduces the recurrence of
\defref{c03-three-orders} verbatim, so the stream equals $\mathit{pre}(v)$ (resp.\
$\mathit{in},\mathit{post}$). Cost is unchanged since the three emit-points are $O(1)$
each.
\end{proofbox}
\begin{insightbox}
Stop memorising three functions. Write \emph{one} walk and move one line. Interviewers
watch for exactly this: whether you relocate the \code{push} or write three near-identical
recursions.
\end{insightbox}

\begin{theorem}{Where the Root Falls}{c03-root-position}
In $\mathit{pre}(\mathit{root})$ the root is the \emph{first} element; in
$\mathit{post}(\mathit{root})$ it is the \emph{last}; in $\mathit{in}(\mathit{root})$ the
root sits at index $\size(L)$, with exactly the left subtree before it and the right
subtree after it.
\end{theorem}
\begin{proofbox}
Immediate from \defref{c03-three-orders}: $\mathit{pre}(v)$ begins with $v$;
$\mathit{post}(v)$ ends with $v$; $\mathit{in}(v)$ is $\mathit{in}(L)$ (length
$\size(L)$), then $v$, then $\mathit{in}(R)$. Hence $v$ is at position
$\size(L)+1$ (index $\size(L)$ zero-based) and the split point of inorder into $L\mid R$
is exactly the root's value.
\end{proofbox}
\begin{corollary}{Reversed Postorder}{c03-rev-post}
$\operatorname{reverse}(\mathit{post})$ is the ``root, right, left'' walk -- the preorder
with children swapped. It is the natural order for \emph{deleting} or \emph{freeing} a
tree iteratively, and for rebuilding bottom-up.
\end{corollary}

\subsection{The Template}

\begin{lstlisting}[caption={One walk, three orders -- move the \code{out.push\_back} line}, label={lst:c03-dfs}]
#include <bits/stdc++.h>
using namespace std;

struct TreeNode {
    int val;
    TreeNode *left, *right;
    TreeNode(int x = 0) : val(x), left(nullptr), right(nullptr) {}
};

vector<int> out;                       // the emitted sequence

void dfs(TreeNode* node) {
    if (!node) return;                 // base case: the empty subtree emits nothing
    out.push_back(node->val);          // <-- HERE = PREORDER  (node, left, right)
    dfs(node->left);
    /* out.push_back(node->val); */    // <-- HERE = INORDER   (left, node, right)
    dfs(node->right);
    /* out.push_back(node->val); */    // <-- HERE = POSTORDER (left, right, node)
}
\end{lstlisting}
\begin{complexitybox}
$\bigO{n}$ time and $\bigO{h}$ stack (worst case $\bigO{n}$ on a skewed tree -- the
\textsf{RE} risk that Section~\ref{sec:c03-iterative} removes). Exactly one of the three
\code{push\_back} lines is active per order; the other two stay commented. The output
\emph{sequence} is what differs, never the walk.
\end{complexitybox}

\subsection{Worked Example}

\begin{example}{The three orders of $R$}{c03-runs}
Applying \coderef{c03-dfs} to $R$ (\figref{c03-running}): reading the contour walk of
\figref{c03-contour} in the three colours gives
\[
\mathit{pre}=1,2,4,5,3,6\quad
\mathit{in}=4,2,5,1,3,6\quad
\mathit{post}=4,5,2,6,3,1.
\]
Check \thmref{c03-root-position}: $1$ is first in \textit{pre}, last in \textit{post}, and
in \textit{in} it sits after exactly its $\size(L)=3$ left-subtree nodes ($4,2,5$) and
before its right subtree ($3,6$). The skewed node $3$ (right child only) shows up as
$\ldots,1,\mathbf{3},6$ in inorder -- $3$ is emitted \emph{before} its right child $6$
because it has no left subtree to stand in the way.
\end{example}
\begin{tracebox}[Preorder emit order on $R$, call-by-call]
\begin{enumerate}[nosep, leftmargin=1.8em]
  \item \code{dfs(1)} emits $1$; descend left
    \begin{enumerate}[nosep,label=(\alph*),leftmargin=1.4em]
      \item \code{dfs(2)} emits $2$; \code{dfs(4)} emits $4$ (leaf), \code{dfs(5)} emits $5$
      \item \code{dfs(3)} emits $3$; \code{dfs(3->left)} returns (null); \code{dfs(6)} emits $6$
    \end{enumerate}
\end{enumerate}
Output stream: $1,2,4,5,3,6$. \hfill ($2\cdot 6 + 1 = 13$ calls, per \lemref{c02-visit-count}.)
\end{tracebox}

\subsection{Variations and Twists}

\begin{itemize}[nosep]
  \item \textbf{N-ary trees}: preorder visits the node then children left-to-right;
        postorder visits children then the node -- \probref{c03-preorder}'s sibling
        problems LC 589/590.                                        \hfill\lv{2}
  \item \textbf{Reverse an order}: emitting at \code{emit\_post} but walking
        (right, left) gives reverse-preorder, the leaf-delete order of
        \corref{c03-rev-post}.                                      \hfill\lv{3}
  \item \textbf{Successor / predecessor questions} are inorder in disguise -- a BST
        returns sorted, so \emph{inorder} is the tool; \secref{c03-iterative}. \hfill\lv{3}
  \item \textbf{Two orders reconstruct the tree} when one is inorder
        (\thmref{c03-root-position}); this is \secref{c03-rebuild}.  \hfill\lv{4}
\end{itemize}

\subsection{Common Pitfalls}\label{sec:c03-dfs-pitfalls}

\begin{warningbox}[Trap 1 -- emitting at the null child, not the node]
Placing the output in the \code{if (!node) \{ out.push\_back(...); \}} branch adds
\textbf{spurious nulls} and visits $n{+}1$ ``things'' instead of $n$. Emit inside the
non-null body, one of the three marked slots.
\end{warningbox}
\begin{lstlisting}[caption={\textcolor{warningred}{WRONG}: pushes on the null base case}, language=C++, basicstyle=\ttfamily\small]
void pre(TreeNode* n, vector<int>& o) {
    if (!n) { o.push_back(-1); return; }   // invents values, off-by-one count
    o.push_back(n->val); pre(n->left, o); pre(n->right, o);
}
\end{lstlisting}
\begin{lstlisting}[caption={\textcolor{successgreen}{RIGHT}: the base case emits nothing}, language=C++, basicstyle=\ttfamily\small]
void pre(TreeNode* n, vector<int>& o) {
    if (!n) return;                        // empty subtree contributes nothing
    o.push_back(n->val); pre(n->left, o); pre(n->right, o);
}
\end{lstlisting}

\begin{warningbox}[Trap 2 -- reusing one \code{out} across test cases]
A \textbf{member/global} \code{vector} is not cleared between the judge's cases, so
case~2's output carries case~1's nodes. Pass the accumulator by reference and clear it,
or return a fresh vector.
\end{warningbox}
\begin{lstlisting}[caption={\textcolor{warningred}{WRONG}: leaked accumulator}, language=C++, basicstyle=\ttfamily\small]
vector<int> out;                       // persists across Solution instances
vector<int> preorderTraversal(TreeNode* r){ dfs(r); return out; }
\end{lstlisting}
\begin{lstlisting}[caption={\textcolor{successgreen}{RIGHT}: local, or cleared on entry}, language=C++, basicstyle=\ttfamily\small]
vector<int> preorderTraversal(TreeNode* r){
    vector<int> out; helper(r, out); return out;   // fresh each call
}
\end{lstlisting}

\begin{warningbox}[Trap 3 -- assuming inorder is ``sorted'' on a \emph{non}-BST]
Inorder yields ascending values \textbf{only} for a binary \emph{search} tree. On the
arbitrary binary tree of \figref{c03-running} it is $4,2,5,1,3,6$ -- not sorted. Never
``find the minimum'' by taking \code{in[0]} on a plain binary tree; that shortcut is a
\chref{c05} property, not a traversal property.
\end{warningbox}
\begin{lstlisting}[caption={\textcolor{warningred}{WRONG}: inorder[0] is NOT the min of a general tree}, language=C++, basicstyle=\ttfamily\small]
vector<int> in = inorder(root); int mn = in[0];      // 4 here, but min is 1!
\end{lstlisting}
\begin{lstlisting}[caption={\textcolor{successgreen}{RIGHT}: the sorted-inorder claim needs a BST}, language=C++, basicstyle=\ttfamily\small]
// valid ONLY when root is a BST (ch.5): in[0] is the minimum.
\end{lstlisting}


\subsection{Problem Lab}

Two \lv{1}--\lv{2} warm-ups solved concisely (they are \coderef{c03-dfs} with the
emit-line in two different slots), then one genuine \lv{4} stretch that never builds a
tree at all -- it turns preorder into a \emph{counting} problem.

\solvedproblem{c03-preorder}{Binary Tree Preorder Traversal}{LeetCode 144}
  {https://leetcode.com/problems/binary-tree-preorder-traversal/}{LC Easy L2}
  {preorder emit-point, recursive}

\begin{problemstatement}[Preorder Traversal]
Given the root of a binary tree, return the nodes' values in \emph{preorder}
(node, left, right). $n \le 100$.
\end{problemstatement}

\solstep{Understanding the Problem}
This is \defref{c03-three-orders}'s $\mathit{pre}$ verbatim: emit $v$, then recurse
left, then right. The only design choice is where the \code{push\_back} sits -- before
both calls (\coderef{c03-dfs}, slot 1).

\solstep{Solution}
Guard $\nil$ (emit nothing), push \code{node->val}, recurse left, recurse right.

\solstep{Implementation}
\begin{lstlisting}[caption={Preorder (LeetCode 144)}, label={lst:c03-preorder}]
class Solution {
    void dfs(TreeNode* node, vector<int>& out) {
        if (!node) return;                 // empty subtree emits nothing
        out.push_back(node->val);          // PRE: before the two calls
        dfs(node->left,  out);
        dfs(node->right, out);
    }
public:
    vector<int> preorderTraversal(TreeNode* root) {
        vector<int> out; dfs(root, out); return out;
    }
};
\end{lstlisting}

\solstep{Trace}
\begin{tracebox}[On $R$]
$1$ (push), go left; $2$ (push), go left; $4$ leaf (push), back; $5$ leaf (push),
back; $2$ done, back; go right of $1$: $3$ (push), left null, $6$ (push).
Output $1,2,4,5,3,6$ \checkmark\ \probref{c03-preorder}.
\end{tracebox}

\solstep{Complexity}
\begin{complexitybox}
Time $\bigO{n}$, space $\bigO{h}$ ($\bigO{\log n}$ balanced, $\bigO{n}$ skewed).
LC's follow-up asks an \emph{iterative} version -- built in \secref{c03-iterative}.
Verdict: \ac.
\end{complexitybox}

\begin{insightbox}[Key takeaway]
Preorder is the ``copy a node on the way in'' order. The \lv{1} skill is relocating one
line; the \lv{2} insight is that a \emph{global} accumulator is a bug (Trap 2) -- thread
\code{out} by reference and clear per case.
\end{insightbox}

\solvedproblem{c03-inorder}{Binary Tree Inorder Traversal}{LeetCode 94}
  {https://leetcode.com/problems/binary-tree-inorder-traversal/}{LC Easy L2}
  {inorder emit-point, recursive}

\begin{problemstatement}[Inorder Traversal]
Return the values of a binary tree in \emph{inorder} (left, node, right). $n \le 100$.
\end{problemstatement}

\solstep{Understanding the Problem}
Move the emit-line to the middle slot of \coderef{c03-dfs}: recurse the left subtree,
\emph{then} push the node, then recurse right. By \thmref{c03-root-position} this lists
every node's left subtree before the node -- the property that later (Chapter~5) makes a
BST come out sorted.

\solstep{Solution}
Recurse left; push \code{node->val}; recurse right; $\nil$ emits nothing.

\solstep{Implementation}
\begin{lstlisting}[caption={Inorder (LeetCode 94)}, label={lst:c03-inorder}]
class Solution {
    void dfs(TreeNode* node, vector<int>& out) {
        if (!node) return;
        dfs(node->left,  out);
        out.push_back(node->val);          // IN: between the two calls
        dfs(node->right, out);
    }
public:
    vector<int> inorderTraversal(TreeNode* root) {
        vector<int> out; dfs(root, out); return out;
    }
};
\end{lstlisting}

\solstep{Trace}
\begin{tracebox}[On $R$]
descend to $4$ (no left) push $4$; up to $2$ push $2$; right $5$ push $5$; up past $2$
to root push $1$; right subtree: $3$ has no left so push $3$, then $6$ push $6$.
Output $4,2,5,1,3,6$ \checkmark\ \probref{c03-inorder}. Notice the skewed node $3$ is
pushed \emph{before} its only (right) child $6$.
\end{tracebox}

\solstep{Complexity}
\begin{complexitybox}
Time $\bigO{n}$, space $\bigO{h}$. Iterative and Morris variants (the latter
$\bigO{1}$ space) are the subject of \secref{c03-iterative}. Verdict: \ac.
\end{complexitybox}

\begin{insightbox}[Key takeaway]
Inorder is the only order that stands \emph{underneath} the node (\figref{c03-contour})
-- it has finished the left subtree but not the right. That asymmetry is why it is the
natural ``sorted read'' and the reconstruction pivot of \secref{c03-rebuild}.
\end{insightbox}

\solvedproblem{c03-verify-pre}{Verify Preorder Serialization of a Binary Tree}{LeetCode 331}
  {https://leetcode.com/problems/verify-preorder-serialization-of-a-binary-tree/}{LC Medium L4}
  {slot / out-degree accounting, single-pass preorder invariant}

\begin{problemstatement}[Verify Preorder Serialization]
A binary tree is written in \emph{preorder} as a comma-separated string: each internal
node is a token, each empty child is \code{'\#'}. Given a string, decide whether it is the
preorder serialization of \emph{exactly one} valid binary tree. Example: \allowbreak
\code{"9,3,4,\#,\#,1,\#,\#,2,\#,\#"} is valid; \allowbreak
\code{"9,\#,\#,1"} is not. Tokens are single digits or \code{'\#'}.
\end{problemstatement}

\begin{attemptgate}[Attempt it yourself first -- \lv{4}, give yourself 25 minutes]
A fair first attempt: try to rebuild the tree from the tokens and check nothing is left
over.
\peek{every node is a hole to fill plus holes it creates -- can you count holes instead
of building anything?}
\nudge{how many child-slots does one node \emph{consume}, and how many does it
\emph{create}? What must be true of this running total after every prefix?}
\fullhint{Track a \code{slots} counter starting at $1$; each token consumes one slot,
a non-\# token then adds two. The string is valid iff slots never go negative and is
$0$ exactly at the end. Full build below.}
\end{attemptgate}

\solstep{Understanding the Problem}
We must judge a \emph{sequence} for tree-legal-ness without a tree. Preorder means
tokens arrive root-first, a whole left subtree, then a right subtree -- so any valid
prefix must have ``opened'' at least as many child positions as it has ``filled.'' The
question is really: \emph{does the string describe a single, exactly-closed binary tree?}

\solstep{Key Observations}
\begin{observation}
Every node, real or \#, occupies exactly one \emph{child-slot}. The whole tree starts
with one open slot (for the root), so \code{slots} begins at $1$.
\end{observation}
\begin{observation}
A \# token is a leaf: it fills its slot and creates none (net $-1$). A non-\# token fills
one slot but opens two children (net $+1$). So \code{slots}-- then, if real, \code{slots
+= 2}.
\end{observation}
\begin{observation}
If \code{slots} ever drops below $0$, we filled a slot that did not exist -- a
too-many-tokens string. The tree is exactly closed iff \code{slots}$=0$ \emph{after the
final} token; $0$ earlier means a complete tree with trailing garbage (invalid).
\end{observation}

\solstep{Think It Through With Me}
\begin{thinkbox}
\deadend{Rebuild the tree with a stack, then verify size.}{Correct but you now maintain
pointers, an index, and a ``did we consume everything'' flag -- the assembly is exactly
\secref{c03-rebuild}, which this problem does \emph{not} ask for. It works, at
$\bigO{n}$ time and stack; we can do it with one integer.}
\deadend{Count the \# and non-\# tokens and demand $\mathit{hash} = \mathit{nonhash}+1$.}
{A full binary tree satisfies $\text{nulls} = \text{internals}+1$, so matching totals are
\emph{necessary} but not \emph{sufficient}: the string \code{"9,\#,\#,\#,1"} has $2$
non-\# ($9,1$) and $3$ \# -- exactly $3 = 2{+}1$, balanced -- yet it is illegal, because the
tree already closes after \code{9,\#,\#} (slots $=0$) and the next \# fills a slot that no
longer exists. Totals ignore \emph{order}.}
\aha{Preorder gives the order for free: keep a running \code{slots} balance that must
stay strictly positive \emph{through} the last token and reach $0$ right after it. A
prefix invariant, not a total count.}
\end{thinkbox}

\solstep{The Approach Ladder}
\begin{center}\renewcommand{\arraystretch}{1.3}
\begin{tabular}{@{}c p{0.235\linewidth} p{0.28\linewidth} l c@{}}
\toprule
\textbf{\#} & \textbf{Approach} & \textbf{Idea} & \textbf{Time} & \textbf{Verdict} \\
\midrule
1 & \lv{2} rebuild + size check & stack-assemble the tree, ensure no leftovers & $O(n)$ & \ac \\
2 & \lv{3} count \# vs non-\# & require $\mathit{hash}=\mathit{nonhash}+1$ & $O(n)$ & \textbf{WA} \\
3 & \lv{4} slot balance scan & prefix \code{slots}$\ge 0$, final $=0$ & $O(n)$, $O(1)$ & \ac \\
\bottomrule
\end{tabular}\end{center}
Rung 2 is the seductive wrong answer (order-blind); rung 3 keeps the preorder order in a
single counter and is the solution.

\solstep{Solution}
Set \code{slots=1}. For each comma-separated token in order: \code{slots--} (fill one
slot); if \code{slots<0} return false (filled a nonexistent slot); if the token is not
\code{'\#'} then \code{slots += 2} (open its two children). After the loop the string is
valid iff \code{slots==0}.

\solstep{Proof of Correctness}
\begin{proofbox}[Invariant on \code{slots}, induction over the token prefix]
\emph{Claim:} after processing any prefix $P$, \code{slots} equals the number of
child-positions of the (partially built) tree described by $P$ that are still awaiting a
token.
\emph{Base:} $P$ empty, \code{slots}$=1$ -- the single slot reserved for the root.
\emph{Step:} the next token fills exactly one awaiting slot (\code{slots--}). If it is
\# it adds none; if it is a node it adds two, so awaiting-slot count changes by $-1$
(resp.\ $+1$), matching the update. Thus the invariant holds.
\emph{Validity:} a slot is filled only when one is open, so if \code{slots} would go
negative the prefix over-fills -- rejected (no valid tree has this prefix). At the end, a
single complete tree has consumed every token into slots and has no open slot left, i.e.\
\code{slots}$=0$; if \code{slots}$>0$ some slots are unfilled (truncated), if a token
arrived when \code{slots}$=0$ the tree already closed and there is extra -- both invalid.
Hence the procedure returns true exactly for valid preorder serialisations.
\end{proofbox}

\solstep{Implementation}
\begin{lstlisting}[caption={Verify Preorder Serialization (LeetCode 331)}, label={lst:c03-verify-pre}]
#include <bits/stdc++.h>
using namespace std;

class Solution {
public:
    bool isValidSerialization(string preorder) {
        int slots = 1;                          // one open slot for the root
        for (int i = 0; i < (int)preorder.size(); ) {
            char c = preorder[i];
            if (--slots < 0) return false;      // filled a slot that was never open
            if (c != '#') slots += 2;           // a real node opens two children
            while (i < (int)preorder.size() && preorder[i] != ',') ++i;  // skip token
            if (i < (int)preorder.size()) ++i;  // step over the comma
        }
        return slots == 0;                      // exactly closed, nothing dangling
    }
};
\end{lstlisting}

\solstep{Trace}
\begin{tracebox}[Running \code{slots} on \code{9,3,4,\#,\#,1,\#,\#,2,\#,\#} (LC's valid sample)]
Start \code{slots}$=1$; each token does $-1$, and each digit then does $+2$. The value
after every token, in order:
$9\!\to\!2$, $3\!\to\!3$, $4\!\to\!4$, $\#\!\to\!3$, $\#\!\to\!2$, $1\!\to\!3$,
$\#\!\to\!2$, $\#\!\to\!1$, $2\!\to\!2$, $\#\!\to\!1$, $\#\!\to\!0$.
The running value is never negative and finishes at $0$ $\Rightarrow$ \textbf{true}.
\end{tracebox}

\solstep{Complexity}
\begin{complexitybox}
Time $\bigO{n}$ (one pass over characters), space $\bigO{1}$ (a single counter). This is
the payoff of rung~3 over rung~1: same time, no stack, no tree.
\end{complexitybox}

\solstep{Verification and Edge Cases}
\begin{center}\renewcommand{\arraystretch}{1.2}
\begin{tabular}{@{}p{0.30\linewidth} p{0.62\linewidth}@{}}
\toprule
\textbf{Case} & \textbf{What the code does} \\
\midrule
Single \code{"\#"} & \code{slots}$1\to0$, end $0$ $\Rightarrow$ true (a lone null tree) \\
\code{"9,\#,\#"} & $1\to2\to1\to0$, true (root with two null children) \\
Truncated \code{"9,\#"} & end \code{slots}$=1\neq0$ $\Rightarrow$ false \\
Over-full \code{"9,\#,\#,\#"} & last \# drives \code{slots}$-1<0$ $\Rightarrow$ false \\
Empty string & loop skipped, \code{slots}$=1\neq0$ $\Rightarrow$ false \\
\bottomrule
\end{tabular}\end{center}

\solstep{Pitfalls and Variants}
\begin{warningbox}
Checking only the \emph{final} \code{slots==0} without the running \code{slots<0} guard
accepts \code{"9,\#,\#,\#"}; and forgetting the ``must be $0$ exactly at the \emph{end}''
check accepts a closed tree plus trailing tokens. Both prefix \emph{and} final conditions
are load-bearing. Variant: the same slot idea verifies preorder of a $k$-ary tree
(each node opens $k$ slots) -- one-argument generalisation.
\end{warningbox}

\begin{insightbox}[Key takeaways -- three things to remember]
\begin{enumerate}[nosep, leftmargin=1.6em]
  \item A traversal \emph{sequence} can often be validated by a running balance instead
        of by rebuilding the structure -- order is information, use it.
  \item ``Necessary total counts'' (rung 2) are a classic trap: a prefix invariant beats a
        global tally whenever order matters.
  \item Slots $=$ (child positions opened) $-$ (filled): the same out-degree bookkeeping
        that makes \thmref{c01-edges-nodes} and handshaking proofs tick.
\end{enumerate}
\end{insightbox}

\begin{hardeningbox}
\begin{itemize}[nosep, leftmargin=1.4em]
  \item \textbf{Adversarial length}: the scan is $\bigO{1}$ space, so it survives strings
        of $10^{5+}$ tokens where a rebuild (rung 1) risks stack overflow on a deep
        skewed tree -- prefer the counter even when rebuilding is available.
  \item \textbf{Multi-digit tokens}: LC restricts values to single digits, but a general
        serialiser must \code{split} on commas and test \code{token == "\#'}. The slot
        logic is unchanged; only the tokenizer grows.
  \item \textbf{Cross-check} with the WA clinic at chapter end: the two most common
        wrong answers here are count-only (rung 2) and final-zero-without-guard.
\end{itemize}
\end{hardeningbox}

\subsection{Practice Problems}

\begin{stuckbox}
\begin{itemize}[nosep, leftmargin=1.4em]
  \item Is your output statement in the \emph{non-null} body, not the base case?
        (Trap 1.)
  \item Are you \emph{clearing} or freshly allocating the output vector each call?
        (Trap 2.)
  \item For a sequence-validation problem: have you found a \emph{running} quantity with
        a prefix invariant, or are you only checking totals? (LC 331 rung 2 vs rung 3.)
  \item On an N-ary node, does your loop keep left-to-right child order?
\end{itemize}
\end{stuckbox}

\begin{practiceset}
  \practice{Binary Tree Postorder Traversal}{LeetCode 145}
    {https://leetcode.com/problems/binary-tree-postorder-traversal/}{LC Easy L2}
    {Same walk, emit in the \emph{third} slot (after both calls); the interesting part is
    LC's iterative follow-up -- see \secref{c03-iterative}.}
  \practice{N-ary Tree Preorder Traversal}{LeetCode 589}
    {https://leetcode.com/problems/n-ary-tree-preorder-traversal/}{LC Easy L2}
    {Push the node, then recurse each child in the \emph{given} (left-to-right) order;
    the children are a vector, so a \code{for} replaces the two fixed calls.}
  \practice{N-ary Tree Postorder Traversal}{LeetCode 590}
    {https://leetcode.com/problems/n-ary-tree-postorder-traversal/}{LC Easy L2}
    {Recurse all children first, then push the node -- postorder's ``children finish
    before me'' invariant, unchanged for many children.}
\end{practiceset}

\begin{recallbox}[10-second refresher -- 3.1 The three depth-first orders]
One contour walk, three read-points: push before both calls $=$ preorder, between $=$
inorder, after $=$ postorder. Preorder leads with the root, postorder ends with it,
inorder splits a node into $L\mid R$. The empty child emits nothing; clear the
accumulator per case.
\selftest{
  \item Write \coderef{c03-dfs} from memory and mark the three emit slots.
  \item On $R$, give all three orders without looking, and say where the root sits in each.
  \item Why is inorder \emph{not} sorted on $R$, and what extra hypothesis would make it so?
  \item State the LC 331 slot invariant and the two conditions (prefix and final) that a
        correct solution must enforce.
  \item What does \emph{reversed postorder} equal, and what is it good for?
}
\end{recallbox}
